\documentclass[../BTL.tex]{subfiles}
\begin{document}

\chapter*{LỜI CẢM ƠN}
\addcontentsline{toc}{chapter}{LỜI CẢM ƠN}

Trong suốt quá trình thực hiện bài tập lớn “Sử dụng thuật toán A* để tìm ra đường đi ngắn nhất trong lưới không gian hai chiều có vật cản”, nhóm chúng em đã nhận được rất nhiều sự giúp đỡ, động viên và hỗ trợ quý báu từ nhiều phía. Đây không chỉ là một bài tập học thuật, mà còn là hành trình rèn luyện tư duy, sự kiên trì và tinh thần làm việc nhóm, để lại cho chúng em nhiều bài học ý nghĩa.

Trước hết, nhóm xin bày tỏ lòng biết ơn sâu sắc tới giảng viên hướng dẫn – ThS. Hoài Thư. Cô không chỉ truyền đạt kiến thức nền tảng một cách rõ ràng, logic mà còn luôn tận tình định hướng, góp ý chi tiết và kịp thời trong suốt quá trình nhóm thực hiện đề tài. Những nhận xét của cô không chỉ giúp nhóm hoàn thiện bài làm mà còn giúp chúng em hiểu sâu hơn về bản chất của thuật toán, cách tư duy trong nghiên cứu và cách tiếp cận một vấn đề một cách khoa học. Sự nghiêm túc và tâm huyết của cô chính là động lực để chúng em cố gắng hoàn thiện bài tập lớn với chất lượng tốt nhất.

Bên cạnh đó, chúng em xin cảm ơn những người bạn, những người đồng đội trong nhóm đã cùng nhau hợp tác, chia sẻ kiến thức, hỗ trợ lẫn nhau và không ngừng cố gắng. Trong quá trình làm việc, chắc chắn không tránh khỏi những lúc bất đồng quan điểm hay gặp khó khăn trong việc triển khai ý tưởng, nhưng chính tinh thần trách nhiệm và sự đoàn kết đã giúp nhóm vượt qua tất cả. Mỗi thành viên đều đã nỗ lực hết mình, đóng góp trí tuệ và thời gian để hoàn thành đề tài một cách trọn vẹn nhất.

Mặc dù đã nỗ lực hết sức, bài tập lớn chắc chắn vẫn còn những hạn chế nhất định. Nhóm rất mong nhận được sự thông cảm và những ý kiến đóng góp quý báu từ giảng viên để có thể tiếp tục hoàn thiện hơn trong tương lai.

Một lần nữa, nhóm xin chân thành cảm ơn tất cả những sự hỗ trợ, đồng hành và tin tưởng đã giúp chúng em hoàn thành bài tập lớn này. Đây sẽ là nền tảng quan trọng để chúng em tiếp tục học tập và phát triển trong lĩnh vực Trí tuệ nhân tạo cũng như trong con đường học thuật phía trước.

\end{document}
